\documentclass[11pt,a4paper]{article}

\usepackage[T1]{fontenc}
\usepackage[utf8]{inputenc}
\usepackage{lmodern}
\usepackage[ngerman]{babel}
\usepackage[a4paper,margin=2.2cm]{geometry}
\usepackage{array}
\usepackage{longtable}
\usepackage{tabularx}
\usepackage{xcolor}
\usepackage{hyperref}

\setlength{\parindent}{0pt}
\setlength{\parskip}{0.5em}
\renewcommand{\arraystretch}{1.35}

\newcommand{\answerline}[1]{\rule{#1}{0.4pt}}
\newcommand{\term}[1]{\fcolorbox{black!20}{blue!6}{\strut #1}}

\title{\textbf{Aufgabenblatt: Kategorien von Systemkomponenten}}
\author{Modul 321 -- Lektion 06}
\date{}

\begin{document}

\maketitle

\textbf{Name:} \answerline{7cm} \hfill \textbf{Datum:} \answerline{3.5cm}

\vspace{0.8em}

\fcolorbox{blue!40}{blue!6}{
  \parbox{\dimexpr\linewidth-2\fboxsep-2\fboxrule-6pt\relax}{
    \textbf{Ziel der Lektion:} Sie können wichtige Kategorien von
    Systemkomponenten benennen, deren Aufgabe erklären und sie in einfachen
    Softwaresystemen sinnvoll einordnen.
  }
}

\section*{Hinweise}
\begin{itemize}
  \item Antworten Sie in ganzen Sätzen oder präzisen Stichworten.
  \item Begründen Sie Zuordnungen kurz und fachlich sauber.
  \item Skizzen und kleine Übersichten sind erlaubt.
\end{itemize}

\section*{1. Begriffe erklären}
Erklären Sie die folgenden Begriffe in \textbf{1 bis 3 eigenen Sätzen}.

\begin{longtable}{|>{\raggedright\arraybackslash}p{0.30\linewidth}|>{\raggedright\arraybackslash}p{0.62\linewidth}|}
\hline
\textbf{Begriff} & \textbf{Ihre Erklärung} \\
\hline
\term{Identitätsverwaltung} & \\[2.1cm] \hline
\term{Ereignisverwaltung} & \\[2.1cm] \hline
\term{Monitoring} & \\[2.1cm] \hline
\term{Datenhaltung} & \\[2.1cm] \hline
\term{Lastenverteilung} & \\[2.1cm] \hline
\end{longtable}

\section*{2. Beispiele zuordnen}
Ordnen Sie jedem Beispiel die \textbf{passendste Komponentenkategorie} zu:

\term{Identitätsverwaltung}, \term{Ereignisverwaltung}, \term{Monitoring},
\term{Datenhaltung}, \term{Lastenverteilung}

\begin{enumerate}
  \item Ein System prüft, ob eine Benutzerin Administratorrechte besitzt.
  \item Eine Meldung \glqq Datei hochgeladen\grqq{} wird an mehrere Dienste weitergegeben.
  \item Ein Dashboard zeigt CPU-Last, Fehlerrate und Antwortzeit.
  \item Kundendaten werden dauerhaft gespeichert und wieder abgefragt.
  \item Eingehende HTTP-Anfragen werden auf drei identische Serverinstanzen verteilt.
  \item Ein Access Token wird nach erfolgreichem Login erstellt.
\end{enumerate}

\section*{3. Komponenten vergleichen}
Vergleichen Sie die folgenden Kategorien.

\begin{tabularx}{\linewidth}{|>{\raggedright\arraybackslash}p{0.28\linewidth}|>{\raggedright\arraybackslash}X|>{\raggedright\arraybackslash}X|}
\hline
\textbf{Frage} & \textbf{Monitoring} & \textbf{Datenhaltung} \\
\hline
Welches Hauptziel hat die Komponente? & & \\[1.1cm] \hline
Welche typischen Informationen verarbeitet sie? & & \\[1.1cm] \hline
Was wäre ein typischer Fehlerfall? & & \\[1.1cm] \hline
\end{tabularx}

\vspace{1em}

\begin{tabularx}{\linewidth}{|>{\raggedright\arraybackslash}p{0.28\linewidth}|>{\raggedright\arraybackslash}X|>{\raggedright\arraybackslash}X|}
\hline
\textbf{Frage} & \textbf{Identitätsverwaltung} & \textbf{Lastenverteilung} \\
\hline
Welche Kernaufgabe hat die Komponente? & & \\[1.1cm] \hline
Welche Probleme löst sie typischerweise? & & \\[1.1cm] \hline
Warum sollte man sie nicht verwechseln? & & \\[1.1cm] \hline
\end{tabularx}

\section*{4. Aussagen beurteilen}
Beantworten Sie mit \textbf{richtig / falsch} und begründen Sie kurz.

\begin{enumerate}
  \item Monitoring speichert in erster Linie Geschäftsdaten wie Bestellungen.
  \item Eine Identitätsverwaltung kann für mehrere Anwendungen gemeinsam genutzt werden.
  \item Lastenverteilung kann helfen, ein System skalierbarer zu machen.
  \item Ereignisverwaltung bedeutet, dass keine Datenhaltung mehr nötig ist.
  \item Eine Datenbank allein ersetzt keine Zugriffssteuerung.
\end{enumerate}

\section*{5. Lernplattform analysieren}
Eine Lernplattform hat folgende Anforderungen:
\begin{itemize}
  \item Benutzer sollen sich anmelden können.
  \item Dateien sollen gespeichert und heruntergeladen werden.
  \item Beim Upload einer Datei sollen andere Teile des Systems informiert werden.
  \item Ausfälle und Fehlerraten sollen sichtbar sein.
  \item Viele gleichzeitige Zugriffe am Morgen sollen abgefangen werden.
\end{itemize}

\begin{enumerate}
  \item Ordnen Sie jeder Anforderung eine passende Komponentenkategorie zu.
  \item Nennen Sie \textbf{eine} Komponente, die aus Sicherheitsgründen besonders sensibel ist.
  \item Nennen Sie \textbf{eine} Komponente, die vor allem für Stabilität und Skalierung wichtig ist.
\end{enumerate}

\section*{6. Fehlende Komponente erkennen}
Ein Team hat eine Webanwendung mit Frontend, API und Datenbank gebaut.
Bei hoher Last wird die Anwendung langsam, und Fehler werden oft erst von Benutzenden gemeldet.

\begin{enumerate}
  \item Welche \textbf{zwei} Komponentenkategorien fehlen offensichtlich oder sind zu schwach ausgeprägt?
  \item Begründen Sie Ihre Antwort.
  \item Welche Probleme könnten dadurch im Betrieb entstehen?
\end{enumerate}

\section*{7. Mini-Entwurf}
Entwerfen Sie für einen einfachen Online-Shop eine grobe Komponentenübersicht.

Funktionen:
\begin{itemize}
  \item Login
  \item Produkte anzeigen
  \item Bestellung auslösen
  \item Benachrichtigung nach Bestellung
  \item Betrieb mit vielen Benutzeranfragen
\end{itemize}

\begin{enumerate}
  \item Nennen Sie mindestens \textbf{vier} passende Komponentenkategorien.
  \item Beschreiben Sie kurz die Aufgabe jeder gewählten Kategorie.
  \item Welche Kategorie wäre für die Benachrichtigung nach einer Bestellung besonders passend?
\end{enumerate}

\section*{8. Reflexion}
Beantworten Sie die Fragen kurz.

\begin{enumerate}
  \item Warum ist es problematisch, wenn eine einzelne Komponente zu viele Verantwortungen übernimmt?
  \item Warum bedeutet \glqq mehr Komponenten\grqq{} nicht automatisch \glqq bessere Architektur\grqq{}?
  \item Was ist bei der Auswahl von Systemkomponenten wichtiger als ein moderner Technologietrend?
\end{enumerate}

\end{document}
